\chapter{Within-subjects experiments}
Blocking increases power and precision by
pairing up similar participants and randomly assigning one of each pair to
each condition. In within-subjects designs, this idea is taken to an extreme.

\mypar{Definition}
  In a \term{within-subjects experiment}, each unit (e.g., participant)
  is assigned to multiple conditions.
\parend

We've already encountered within-school and within-class designs
in the previous chapter;
within-subjects designs extend the same logic to the individual unit.

Typically, the units in a within-subjects design are assigned to all
of the experiment's conditions, though it is possible to devise
within-subjects experiments in which each unit is assigned to
only a subset of the experiment's conditions
\citep[see][Section 7.3]{Senn2002}.
We will only discuss within-subjects experiments in which 
each unit is assigned to all conditions, however.

\section{Advantages and drawbacks}
The main advantage of within-subjects designs is 
that they tend to have greater statistical precision and power than an otherwise
comparable experiment in which each unit is assigned to only one condition.
Indeed, a study's statistical precision depends on
 (a) the amount of data and
 (b) the variability in the data. 
Within-subjects designs easily account for a major source of variability in the data:
pre-existing differences between participants.
 
How much more precise a within-subjects experiment is than a
between-subjects experiment varies from case to case.
\citet{Quene2010} 
estimates that within-subjects designs have the statistical precision
of between-subjects designs with four times as many participants.
The precise factor depends on the extent to which the participants'
performance in one condition correlates with their performance in the
other condition:
The stronger this correlation, the greater the added value of a
within-subjects experiment.
But even if you can't quantify this added value: 
Within-subjects designs typically offer greater statistical precision.
 
Within-subjects designs, however, also suffer from a few potential drawbacks.
First, in applied settings, you typically want the study 
to mimic the context in which its findings are to be implemented.
But in such a context, the units (e.g., pupils) won't be exposed to several 
conditions (e.g., learning methods) but rather to just one.
Hence, the within-subjects design may be criticised for a lack of ecological validity.

A second potential drawback is that,
when participants are tested in several conditions,
it's possible that they learn something in one condition
that affects their performance in the other condition.
This phenomenon is called a \term{carry-over effect}.
It's also possible that their performance in the last condition differs
from their performance in the first condition because they've grown
accustomed to the setting or because they've grown tired of being tested.
These phenomena are referred to as
\term{learning effects}, \term{fatigue effects},
and more generally, \term{period effects}.

\section{Minimising period effects}
The first technique to prevent period effects 
from biasing the results is to vary 
the sequence in which the conditions are presented from unit to unit.
In \term{complete counterbalancing}, \emph{all} possible sequences are 
taken into account. If you have two within-subjects conditions,
half of the participants first complete condition $A$ and then $B$,
and the other half first complete condition $B$ and then $A$.
If you have three conditions, there are 3! = 6 possible orders;
you can then randomly assign one sixth of the participants to 
first complete A, then B, then C;
one sixth completes A, then C, then B, etc.:
 
\begin{table}[h]
\centering
\caption{Complete counterbalancing for a within-subjects experiment with three conditions.}
\label{tab:counter}
\begin{tabular}{cccc}
\toprule
Sequence & Period 1 & Period 2 & Period 2 \\
\midrule
\textsc{i}  & A & B & C \\
\textsc{ii} & A & C & B \\
\textsc{iii}&  B & A & C \\
\textsc{iv} & B & C & A \\
\textsc{v}  & C & A & B \\
\textsc{vi} & C & B & A \\
\bottomrule
\end{tabular}
\end{table}

\medskip

If a within-subjects experiment has lots of conditions,
complete counterbalancing is impractical. For four conditions,
for instance, there are already 4! = 24 possible orders---we may not
even have that many participants! 
The second technique, the \term{Latin square design}, lends itself
to such cases.
Latin squares are arrangements of symbols in a grid in which
each of the symbols used occurs exactly once in each row
and exactly once in each column,
and they are called so because the symbols are usually letters from
the Latin alphabet.
The grid in Table \ref{tab:latquad} is a Latin square
of size 4---one of the 576 possible arrangements of the symbols A, B, C, and D
that form a Latin square.
 
\begin{table}[h]
\centering
\caption{A Latin square for a within-subjects experiment with four conditions.}
\label{tab:latquad}
\begin{tabular}{ccccc}
\toprule
Sequence & Period 1 & Period 2 & Period 3 & Period 4 \\
\midrule
\textsc{i}  & A & B & C & D \\
\textsc{ii} & B & C & D & A \\
\textsc{iii}& C & D & A & B \\
\textsc{iv} & D & A & B & C \\
\bottomrule
\end{tabular}
\end{table}

\mypar{Exercise}
There are four different Latin squares of size 4 in which the first column
and the first row are both in natural order, i.e., ABCD.
Such a Latin square is said to be in reduced form.

Table \ref{tab:latquad} shows one Latin square of size 4 in reduced form.
Construct the other three.
\parend

When adopting a Latin square design,
you first pick a Latin square of the appropriate size in reduced form.
It's reasonable to pick this Latin square at random,
but you can also, for instance, pick a square 
in which each sequence in the square (e.g., ABCD) also occurs in reverse
order (e.g., DCBA).
You then randomly assign the participants to the sequences defined by
the square's rows.
Finally, you assign the conditions randomly to the symbols used.

For instance, let's say you picked the Latin square in 
Table \ref{tab:latquad} for your study.
You'd then randomly assign one quarter of the participants to
the condition (or stimulus set) order ABCD (first row), 
one quarter to BCDA (second row),
one quarter to CDAB (third row) and one quarter to DABC (fourth row).\marginnote{Don't worry
if the different sequences aren't represented \emph{exactly} equally frequently, though.
Whatever you do, do \emph{not} throw away data just to make each sequence occur equally often!
This comment also applied to completely counterbalanced designs.}
The conditions (or stimulus sets) are randomly
assigned to one of the letters, too.
The participants assigned to the first row complete the conditions (or stimulus sets)
in the order ABCD;
those assigned to the second row complete them in the order BCDA, etc.

\mypar{Tip}
  If your within-subjects experiment features three conditions,
  complete counterbalancing is a good default strategy.
  If it features four or more conditions,
  a Latin square design keeps things reasonably simple.
\parend

Counterbalancing and Latin squares do \emph{not} negate carry-over effects.
Whether carry-over effects represent an acute danger to the study's
validity is a subject-matter decision (not a statistical one)
that needs to be settled \emph{beforehand} on a case-by-case basis.
Within-subjects designs are not suited for contexts in which carry-over effects
are likely to be substantial.
That said, the possible danger of carry-over effects quite often isn't
large enough to offset the certain gain in statistical precision.

\begin{framed}
If you have a genuine choice between a between-subjects
and a within-subjects design for your own research, pick the 
within-subjects design. (Unless, of course, you have an excellent
reason not to do so.)
\end{framed}

\mypar{Example}
  In \citet{Vanhove2018}, I examined the influence of metalinguistic knowledge
  about some structure in the L1 on the participants' intuitions about 
  the corresponding structure in a foreign language.
  Metalinguistic knowledge was manipulated experimentally.
  In the control condition (call it $A$), the participants were 
  given correct but irrelevant information about the L1's morphosyntax.
  In the second condition ($B$), they were given correct, relevant information.
  In the third condition ($C$), they were given correct, relevant information
  and they were additionally taught an algorithm for identifying the structure in question.
  
  It's quite clear that this study couldn't be realised in a within-subjects design:
  Once you give the participants the relevant information (conditions $B$ and $C$), 
  you can't then expect them to suppress this information as they are tested in condition $A$.
  Similarly, once you teach the participants an algorithm for identifying a structure (condition $C$),
  you can't then expect them not to apply it in the other conditions.
  In view of such carry-over effects,
  I opted for a between-subjects design.
\parend

\mypar[Potential outcomes in within-subjects designs]{Remark}
  We introduced the notion of potential outcomes when discussing 
  true (between-subjects) experiments in Chapter \ref{ch:randomisation}.
  If each unit participates in both conditions of a two-condition within-subjects
  experiment, it may seem as all potential outcomes are, in fact, known.
  But this is not the case.
  For a unit assigned to the $AB$ order, the known outcomes are how they
  performed in the $A$ condition \emph{if this condition is presented first}
  and how they performed in the $B$ condition \emph{if this condition is presented second}.
  But there are still the unknown potential outcomes of how they would have performed
  in the $A$ condition if it had been presented second and in the $B$ condition
  if it had been presented first.
\parend

The preceding discussion assumes that the experiment's conditions
have to appear one after the other.
Sometimes, however, the conditions can be interleaved.
In particular, this is the case in psycholinguistic studies in which
the participants need to react to several stimuli per condition
(e.g., 12 sentences per condition).
In such settings, each participant is, in a sense, tested several times
in each condition, and the order of presentation can be randomised
for each participant (e.g., ABAABBBAB etc.).

Similarly, it is sometimes desirable to randomise the order of the stimuli
for each participant.
For instance, let's say that we want to run a study
in which we ask the participants to describe 50 pictures,
which we show to them successively. 
If we also want to draw comparisons between the descriptions 
of the different pictures, 
then we cannot just present them in the same order to all participants: 
order effects could plausibly
bias the results.
But 50! = 3 $\cdot$ 10\textsuperscript{64} is an astronomical number, ruling out
complete counterbalancing.
Even a Latin square of size 50 seems impractical.
One possible solution is to present the images in a new random order for each
participant.
The drawback of doing this is that perhaps image 3 occurs 
much more often at the start than at the end of the data collection,
but at least this imbalance isn't built into the design of the study.

Of course, it's possible that we just accept such `bias' if we aren't interested
in differences between the images, but just in differences between
the participants or between the conditions. 
If that's the case, randomising the stimulus order for each participant is superfluous.

\mypar{Exercise}
From \citet{Ludke2014}:
 \begin{quote}
Participants were randomly assigned to one of three learning conditions:
speaking, rhythmic speaking, and singing. The participants heard 20 paired-associate
phrases in English and an unfamiliar language (Hungarian) (\dots).
(\dots) The 15-min learning period was followed by a series of five different production,
recall, recognition, and vocabulary tests for the English--Hungarian pairs.
\end{quote}
Re-design this between-subjects experiment as a within-subjects experiment.
What would this description look like? For the time being, ignore the \textit{rhythmic speaking} condition.

In your design, you are constrained by the resources Ludke et al. had. This means that
you cannot introduce stimuli and tests that were not already used by Ludke et al., and
that you have to work with at most 30 men and 30 women as participants.
\parend

\mypar{Exercise}
As in the previous exercise, but with all three conditions.
\parend

\mypar{Exercise}
A researcher wants to study if foreign-language vocabulary is best learnt
when rehearsing the words one day (`short spacing')
or one week after first studying them (`long spacing').
He designs a within-subjects experiment in which fifty participants
learn twenty items across the two conditions.
Specifically, ten of the items are selected as `short spacing' items,
and ten as `long spacing' items.
On the first day, all participants study all twenty items.
On the next day, they rehearse the `short spacing' items;
on the eighth day, they rehearse the `long spacing' items.
The next day, so on the ninth day, they are tested on all items;
a few weeks later, they are tested again on all items.

\begin{enumerate}[(a)]
  \item Can this study tell us if these participants learn better with short spacing or with long spacing?
        Briefly justify your answer.
        
  \item Suggest an improvement to this design. \parend
\end{enumerate}

\section{*Analysing data from within-subjects designs with two conditions}
\subsection{Two conditions with complete counterbalancing}
Consider a within-subjects experiment that has two conditions and uses
complete counterbalancing.
We will consider condition $A$ to be the intervention (or treatment) condition
and $B$ the control condition, but nothing hinges on these labels.
Let's break down the systematic factors that contribute 
to the measurements in both period for both sequences, 
see Table \ref{tab:analysiswithin}.
  
\begin{table}[h]
  \centering
  \caption{Systematic factors contributing to the measurements in a within-subjects experiment with two conditions. The meaning of the Greek symbols is explained in the running text.}
  \label{tab:analysiswithin}
  \begin{tabular}{ccc}
    \toprule 
    Sequence& First period      & Second period \\
    \midrule
    AB      & $\beta + \tau$    & $\beta + \pi + \kappa$ \\
    BA      & $\beta$           & $\beta + \tau + \pi$ \\
    \bottomrule
  \end{tabular}
\end{table}

\medskip

\begin{itemize}
  \item We assume that there is some baseline $\beta$ common to all measurements.
        This common baseline is of no further interest.
  
  \item There may be a differential treatment effect $\tau$
        that contributes to the measurements in condition $A$ 
        (i.e., the first period for sequence $AB$ and the second
        for sequence $BA$), but not to those in condition $B$.
        It is this different treatment effect we're interested in.
      
  \item There may be a differential period effect $\pi$ that contributes to 
        the measurements in the second period but not to those in the first.
        
  \item There may be a differential carryover effect $\kappa$ 
        that affects the measurements in condition $B$ 
        but only if $B$ follows $A$.
\end{itemize}

In any analysis of these data, 
we need to assume that there are no carryover effects, that is, $\kappa = 0$:
If $\kappa \neq 0$, 
the analysis presented here is no longer appropriate,
and a within-subjects design may not have been a good choice.
Under the assumption that $\kappa = 0$, 
we may estimate the value of $\tau$ by first computing,
for each participant, the period difference, that is, 
the difference between their first and their second measurement.
For the participants assigned to the AB sequence, and ignoring any non-systematic effects, 
we obtain
\[
  d_{AB} = (\beta + \tau) - (\beta + \pi + \kappa) = \tau - \pi - \kappa = \tau - \pi,
\]
since $\kappa = 0$ by assumption.
For the participants assigned to the BA sequence, we similarly obtain
\[
  d_{BA} = \beta - (\beta + \tau + \pi) = - \tau - \pi.
\]
Observe that
\[
  \frac{d_{AB} - d_{BA}}{2} = \frac{(\tau - \pi) - (- \tau - \pi)}{2} = \tau.
\]
The consequence of this is that, under the assumption of no carryover effects,
the treatment effect $\tau$ can be estimated as half the mean difference in the period
differences between the two orders.
A randomisation test or some analytical approximation thereof 
can be used to obtain a $p$-value for the null hypothesis
that $\tau = 0$.

\mypar{Remark}\label{rem:averageofaverages}
  Alternatively, we may, for each participant, compute the difference
  between their outcomes in the $A$ and $B$ conditions.
  For the participants assigned to the AB sequence,
  we then again obtain $\tau - \pi$;
  for those assigned to the BA sequence,
  we obtain $\tau + \pi$.
  Hence, we may also first compute the averages of the $A - B$ differences
  for both sequences, and then compute the mean of these two averages.
  This procedure is mathematically equivalent to the one described in the main text,
  but the randomisation test for it is slightly trickier to implement correctly.
\parend

\mypar[Generalisation to multiple conditions]{Remark}
  The advantage of the `average of averages' approach outlined in
  Remark \ref{rem:averageofaverages} is that it can straightforwardly
  be extended to designs with multiple conditions.
  For instance, if we want to compare the $A$ and $B$ conditions in a
  completely counterbalanced design with three conditions,
  we can compute the average $A - B$ difference for each of the six
  sequences, and then average these six differences.
\parend

\mypar{Remark}
In practice, you'll encounter different approaches to analysing within-subjects 
designs when reading social science studies.
Some of these are, in fact, equivalent to the approach outlined above,
but others tacitly assume that $\pi = 0$ and therefore may yield different
results.
\parend

\mypar{Remark}
  If we aren't willing to exclude the possibility that there is 
  a carryover effect (i.e., that $\kappa \neq 0$),
  we can restrict the analysis to the first measurement
  and treat the data as stemming from a true (between-subjects) experiment.
  However, by doing so, we lose half of the data as well as
  any improvement in statistical precision that we sought to obtain
  by adopting a within-subjects design.
\parend

\subsection{Interleaved conditions}
Consider again a within-subjects experiment with two conditions, $A$ and $B$,
in which $n$ participants are shown a fixed list of $m$ words $w_1, \dots, w_m$.
For each participant, half of the words are shown in condition $A$,
and half are shown in condition $B$;
the presentation software assigns the words to the conditions randomly
and separately for each participant.
The participants react to each word, and their response is expressed numerically
in some fashion (e.g., accuracy, speed, \dots).

For a design like this, there is no real sense in which one condition
follows another condition for each participant.
A sensible way to analyse these data is to compute, for each participant,
the difference score between their average response to condition $A$
and their average response to condition $B$, resulting in difference
scores $d_1, \dots, d_n$.
We then compute some average of these differences, $\overline d$
(e.g., the mean or the median).
Under the null hypothesis of no difference,
each $d_i$ value could just as likely have had the opposite sign.
To generate the distribution of $\overline d$ under the null hypothesis,
we can randomly flip the signs of the observed $d_i$ values
and recompute their average.

Exhaustive sign-flipping would generate the full distribution of $\overline d$
under the null hypothesis, but it involves generating $2^n$ $\overline d$ values.
A Monte Carlo version of this procedure would also work.
Once the distribution of $\overline d$ has been generated or approximated using this sign-flipping,
$p$-values can be computed as usual.
An analytical short-cut to this procedure is the paired $t$-test on the condition averages per participant.
  
\mypar[Within-school/class designs]{Remark}
  In within-school or within-class designs (see Figure \vref{fig:withinschool}),
  the same procedure can be applied: For each school (or class), compute
  the average score obtained by the pupils in the intervention and the average
  score obtained by the pupils in the control condition.
  Then analyse these averages as described above.
\parend

\section{*Further reading}
Latin-square designs are also used in studies other than 
within-subject experiments, for instance as a technique for blocking.
See \citet{Richardson2018} for an overview and some finer points that weren't discussed here; 
his article is geared towards educational researchers.

\citet{Senn2002} discusses the design and analysis of within-subjects experiments
(known as cross-over trials in medical circles)
with an audience of medical researchers and applied statisticians in mind.

\section{*Examples in R}
See the file \texttt{within.html} in the \texttt{tutorials} directory
on \url{https://github.com/janhove/QuantitativeMethodology}.
